\subsection{Метод реологических моделей}

\begin{warn}
Подзаголовок в тетради: «(принцип объективности материальный)» — запись
неразборчива.
\end{warn}

Цель — из уже имеющихся материалов собрать новые модели. Каждый элемент
$\alpha = 1, 2, 3, \dots, N$ нагружен девиатором $\Sdev_\alpha$ и деформируется
с девиатором $\dev_\alpha$; элементом может быть пружина, демпфер или элемент
сухого трения.

\inputfig{figures/fig-rheo-block}

Комбинируем элементы, чтобы сделать новый материал.

\subsubsection*{1) Параллельное соединение}

\inputfig{figures/fig-rheo-connections}

Поперечины движутся поступательно, поэтому деформации одинаковы —
это \textbf{аксиома параллельного соединения}:
\[
  \dev_\alpha = \dev_\beta ,
  \qquad
  \Sdev = \Sdev_\alpha + \Sdev_\beta ,
\]
\[
  \psi_\kappa = \psi_{\kappa\alpha} + \psi_{\kappa\beta}
  \quad(\text{плотность свободной энергии}),
  \qquad
  s_\kappa = s_{\kappa\alpha} + s_{\kappa\beta}
  \quad(\text{энтропия}).
\]

\subsubsection*{2) Последовательное соединение}

\[
  \Sdev = \Sdev_\alpha = \Sdev_\beta ,
  \qquad
  \dev = \dev_\alpha + \dev_\beta ,
\]
\[
  \psi_\kappa = \psi_{\kappa\alpha} + \psi_{\kappa\beta} ,
  \qquad
  s_\kappa = s_{\kappa\alpha} + s_{\kappa\beta} .
\]

Эти аксиомы позволяют гарантировать, что новые материалы — «хорошие»
(удовлетворяют принципам).

\begin{warn}
Индекс у $\psi$ и $s$ в тетради написан кириллической «к» — здесь он набран как
$\kappa$. Что именно он помечает (конфигурацию? формоизменение?), на лекции не
расшифровано.
\end{warn}

\subsubsection*{Проверка II закона термодинамики}

Неравенство Клаузиуса–Дюгема для формоизменения:
\[
  \Sdev \ddotp \dt{\dev} - \rho_0 \dt\psi_\kappa - \rho_0 s_\kappa \dt T \ge 0 .
\]
Оно должно выполняться для новых материалов. Проверим его для параллельного
соединения:
\[
  \Sdev \ddotp \dt{\dev}
  - \rho_0\bigl(\dt\psi_{\kappa\alpha} + \dt\psi_{\kappa\beta}\bigr)
  - \rho_0\bigl(s_{\kappa\alpha} + s_{\kappa\beta}\bigr)\dt T \ge 0 .
\]
Материалы $\alpha$ и $\beta$ — хорошие, мы это знаем:
\[
  \begin{cases}
    \Sdev_\alpha \ddotp \dt{\dev}_\alpha - \rho_0 \dt\psi_{\kappa\alpha}
      - \rho_0 s_{\kappa\alpha}\dt T \ge 0, \\[0.4em]
    \Sdev_\beta \ddotp \dt{\dev}_\beta - \rho_0 \dt\psi_{\kappa\beta}
      - \rho_0 s_{\kappa\beta}\dt T \ge 0 .
  \end{cases}
\]
Тогда сложим (учитывая $\dev_\alpha = \dev_\beta = \dev$):
\[
  \underbrace{\bigl(\Sdev_\alpha + \Sdev_\beta\bigr)}_{\Sdev} \ddotp \dt{\dev}
  - \rho_0\bigl(\dt\psi_{\kappa\alpha} + \dt\psi_{\kappa\beta}\bigr)
  - \rho_0\bigl(s_{\kappa\alpha} + s_{\kappa\beta}\bigr)\dt T \ge 0 .
\]
То есть параллельное соединение — хороший материал (удовлетворяет II закону
термодинамики).
